\chapter{A Catalogue of \texorpdfstring{$\Phi$}{Phi}: Seven Families, Twenty-Eight Types}
\label{ch:catalog}
%============================================================

\section{How the catalogue is built}

Within the space generated by the six degrees of freedom of Section~\ref{sec:dof}, the
points and regions that recur are organized into seven families and listed. For each type
the sign of the credit position $\kappa$, the character of settlement, and representative
examples are given. The catalogue claims no exhaustiveness; it is a working classification
used to select objects of study.

\section{Family 1: one-off exchange}

Forms in which one $\delta$ corresponds to one $\pi$.

\begin{center}
\begin{tabularx}{\textwidth}{@{}llXX@{}}
\toprule
No. & Type & $\kappa$ / character of settlement & Examples \\
\midrule
1-1 & immediate exchange & $\approx 0$ / simultaneous with $\delta$ & retail, restaurants, vending \\
1-2 & one-off advance & $<0$ / precedes $\delta$ & made-to-order, weddings, tickets, crowdfunding \\
1-3 & one-off arrears & $>0$ / follows $\delta$ & B2B trade credit, invoiced sales \\
1-4 & instalments & $>0$ (long) / one $\delta$ split into many $\pi$ & hire purchase, BNPL, mortgages \\
\bottomrule
\end{tabularx}
\captionof{table}{Family 1: one-off exchange.}
\label{tab:family-1-catalog}
\end{center}

Only 1-4 is qualitatively different in structure. The delivery $\delta$ occurs at a point in
time while $\pi$ is spread over a long horizon; it commercializes the difference in discount
factors of Section~\ref{sec:objective} on its own. Its surplus is weighted towards
$\phi_{\mathrm{barg}}$ and $\phi_{\mathrm{cog}}$.

\section{Family 2: continuing access}

Forms in which a right $\dbar$ is sold and $\delta$ is exercised against it. The principal
home of cognitive surplus.

\begin{center}
\begin{tabularx}{\textwidth}{@{}llXX@{}}
\toprule
No. & Type & $\kappa$ / character of settlement & Examples \\
\midrule
2-1 & flat access & annual $<0$ / monthly in arrears $>0$, $\pi=F$ & SaaS, membership retail, facility use \\
2-2 & usage & $\approx 0$ / $\pi = p\,\delta$ & electricity, cloud, telecoms \\
2-3 & two-part tariff & mixed / $\pi = F + p\,\delta$ & mobile, dues plus goods, booking discounts \\
2-4 & prepaid & $<0$ / $\pi$ first, timing of $\delta$ at the customer's discretion & gift cards, ticket books, transit cards \\
2-5 & options and warranties & $<0$ / $\pi$ first, $\delta$ only on exercise & extended warranties, cancellable bookings, credit lines \\
\bottomrule
\end{tabularx}
\captionof{table}{Family 2: continuing access.}
\label{tab:family-2-catalog}
\end{center}

2-1 divides in two according to whether a capacity constraint is present (the proposition of
Section~\ref{sec:breakage}). 2-4 and 2-5 are formally sales of options, leaving the operator
short.

\section{Family 3: renting an asset over time}

\begin{center}
\begin{tabularx}{\textwidth}{@{}llXX@{}}
\toprule
No. & Type & $\kappa$ / character of settlement & Examples \\
\midrule
3-1 & lease and rental & the asset comes first / proportional to time & equipment leasing, property rental \\
3-2 & shared asset & the asset comes first / tied to utilization & lodging, car sharing, coworking \\
3-3 & transfer plus long collection & $>0$ (large) / $\pi$ continues after title passes & vendor finance, equipment finance \\
\bottomrule
\end{tabularx}
\captionof{table}{Family 3: renting an asset over time.}
\label{tab:family-3-catalog}
\end{center}

What characterizes this family is that the dominant term is not $\kappa$ but the $A$
of~\eqref{eq:cash}. In 3-2 a capacity constraint is explicitly present, so the utilization
rate decides feasibility.

\section{Family 4: recovery through a complement}

Forms that split $\delta$ and deliberately skew $\pi$.

\begin{center}
\begin{tabularx}{\textwidth}{@{}llXX@{}}
\toprule
No. & Type & $\kappa$ / character of settlement & Examples \\
\midrule
4-1 & lock-in consumables & early $>0$ / $\delta_1$ below cost, recovered on $\delta_2$ & printers, coffee machines, games consoles \\
4-2 & freemium & early $>0$ / $\pi=0$ for some, recovered from converts & storage, business chat, mobile games \\
4-3 & hardware subsidy plus finance & $>0$ / device discounted, recovered through recurring charges & mobile carriers, auto sales finance \\
\bottomrule
\end{tabularx}
\captionof{table}{Family 4: recovery through a complement.}
\label{tab:family-4-catalog}
\end{center}

By definition this family deliberately creates $\kappa>0$ early on. If the payback period on
customer acquisition lengthens, the whole family becomes dependent on capital markets.

\section{Family 5: outcome- and state-contingent}

Forms in which $\pi$ is measurable with respect to $\omega$. The arguments of
Chapter~\ref{sec:info} bear directly here.

\begin{center}
\begin{tabularx}{\textwidth}{@{}llXX@{}}
\toprule
No. & Type & $\kappa$ / character of settlement & Examples \\
\midrule
5-1 & success fee & $>0$ (large) / $\pi=h(\text{outcome})$ & recruitment, M\&A advisory \\
5-2 & revenue share & $>0$ / $\pi = $ customer revenue $\times$ rate & franchising, app distribution, agencies \\
5-3 & cost plus & intermediate / $\pi = $ cost $+$ margin & some construction, consulting, public procurement \\
5-4 & underwriting & $<0$ / $\pi$ taken first, $\delta$ depends on $\omega$ & life and non-life insurance, guarantees \\
5-5 & IP licensing & mixed / royalties, zero marginal cost of $\delta$ & semiconductor IP, character licensing, patents \\
\bottomrule
\end{tabularx}
\captionof{table}{Family 5: outcome- and state-contingent.}
\label{tab:family-5-catalog}
\end{center}

5-1 and 5-2 transfer risk because effort is unobservable; 5-3 transfers it in the opposite
direction because the outcome is uncertain (the remark in Chapter~\ref{sec:info}).
Only 5-4 has a mirrored time structure within the family.

\section{Family 6: intermediation}

Forms that stand between two other parties' transaction.

\begin{center}
\begin{tabularx}{\textwidth}{@{}llXX@{}}
\toprule
No. & Type & $\kappa$ / character of settlement & Examples \\
\midrule
6-1 & transaction fee & $\approx 0$ / volume $\times$ rate & securities, property brokerage, marketplaces \\
6-2 & escrow & $<0$ (other people's funds) / held then remitted & payment processing, lodging platforms, travel \\
6-3 & spread & $>0$ (inventory) / the gap between bid and ask & wholesale, currency exchange, market making \\
6-4 & charging for opportunity & $<0$ / charged for listing or contact, not the transaction & job advertising, search advertising \\
\bottomrule
\end{tabularx}
\captionof{table}{Family 6: intermediation.}
\label{tab:family-6-catalog}
\end{center}

The $\kappa$ of 6-2 is not the operator's own money. It occupies the same place
in~\eqref{eq:cash} as other negative $\kappa$ but differs in that it drains rapidly when
volume falls. The difference between 6-1 and 6-3 is whether inventory risk is held, which is
again the transfer question of Chapter~\ref{sec:info}.

\section{Family 7: separating beneficiary from payer}

The operation of moving the index $i$ itself.

\begin{center}
\begin{tabularx}{\textwidth}{@{}llXX@{}}
\toprule
No. & Type & $\kappa$ / character of settlement & Examples \\
\midrule
7-1 & advertising & mixed / $\pi=0$ for the beneficiary, a third party pays & broadcast, search, social \\
7-2 & cross-subsidy & mixed / one side free, recovered from the other & single-sided charging in two-sided markets \\
7-3 & public payment & $>0$ / an insurer or government pays & healthcare, long-term care, education \\
7-4 & donation and support & $<0$ / no explicit consideration & support-based charging, tipping, non-profits \\
\bottomrule
\end{tabularx}
\captionof{table}{Family 7: separating beneficiary from payer.}
\label{tab:family-7-catalog}
\end{center}

7-3 is special: because the payer sets the price, the design freedom of $\Phi$ is absorbed
into the institutional layer. A revision of the official price changes $\kappa$ and $m$
exogenously and at the same time.

\section{Assigning an observation to a type}
\label{sec:type-assignment}

To use the catalogue empirically, a rule is needed for assigning an observed $\Phi$ to a
type. A rule decided after the fact allows the assignment to be chosen to fit the
conclusion. The rule is fixed below.

The information needed for assignment is the degrees of freedom of Section~\ref{sec:dof},
which coincide with the variables listed in Chapter~\ref{ch:next} as those to be recorded
for each $\Phi$.

\textbf{The conditions defining the families are not mutually exclusive}, so the order of
the tests is fixed.

\begin{center}
\begin{tabularx}{\textwidth}{@{}llX@{}}
\toprule
Order & Test & If it applies \\
\midrule
1 & Do beneficiary and payer differ (degree of freedom (3))? & family 7 \\
2 & Does it stand between two other parties' transaction? & family 6 \\
3 & Is $\pi$ measurable with respect to $\omega$ (degree of freedom (4))? & family 5 \\
4 & Is delivery split and $\pi$ deliberately skewed? & family 4 \\
5 & Does an asset come first, recovered in proportion to time or utilization? & family 3 \\
6 & Is there repetition (degree of freedom (6))? & family 2; otherwise family 1 \\
\bottomrule
\end{tabularx}
\captionof{table}{Order of the family tests. Applied from the top, the first family that applies is taken.}
\label{tab:family-assignment-order}
\end{center}

Once the family is fixed, the number within it is specified by degree of freedom (1)
(timing), (2) (form of dependence) and (5) (separation of right from exercise).

\begin{remark}[Some assignments depend on the order]
\label{rem:assignment-order-dependence}
There exist $\Phi$ whose assignment changes if the order in
Table~\ref{tab:family-assignment-order} is changed.

An app-distribution fee applies a rate to the volume distributed, which fits 6-1
(transaction fee), and moves with the seller's revenue, which fits 5-2 (revenue share).
Under the order above, family 6 is tested first.

\textbf{There is no deductive ground for the order.} It follows a policy of testing the
composition of parties before the form of settlement, and is a convention. Being a
convention is precisely why it must be fixed in advance.
\end{remark}

\begin{remark}[Handling composite contracts]
\label{rem:composite-assignment}
As Section~\ref{sec:composite} notes, actual contracts are often composites of several types.
In that case $\Phi$ is first decomposed into components and
Table~\ref{tab:family-assignment-order} applied to each.

The unit of decomposition is the range over which one correspondence between delivery and
settlement closes. Selling a device and supplying consumables are separate correspondences,
so they split into two components, assigned to 1-1 and 2-2 respectively and recorded as the
composite 4-1.
\end{remark}

\section{Limits of the catalogue}

The twenty-eight types above are not exhaustive. Of three omissions considered, two can be
described with existing types.

\subsection{The \texorpdfstring{$\Phi$}{Phi} of labour: no family is added}
\label{sec:labor-phi}

An employment contract is also a map from delivery to settlement, but the direction is
reversed: the worker supplies $\delta$ and the firm pays $\pi$. Because the sign convention
for $X$ in Chapter~\ref{sec:phi} makes inflows to the firm positive, all twenty-eight types
presume the firm is the side supplying $\delta$.

But \textbf{seen from the worker's side, the worker is an operator supplying $\delta$}.
The structure is the same as the solo business of Chapter~\ref{ch:solo}, and the catalogue
applies unchanged.

\begin{center}
\begin{tabularx}{\textwidth}{@{}lXl@{}}
\toprule
Contract form & Type & $\kappa_L$ \\
\midrule
Monthly salary in arrears & repeated 1-3, one-off arrears & $>0$ (the worker extends credit) \\
Daily or weekly pay & 1-1, immediate exchange & $\approx 0$ \\
Advances and signing payments & 1-2, one-off advance & $<0$ \\
Bonuses & 5-1, success fee & $>0$ \\
Equity compensation & 5-2, revenue share & $>0$ (extreme) \\
Retirement payments & 1-4, instalments & $>0$ (very long) \\
\bottomrule
\end{tabularx}
\captionof{table}{Contract forms seen from the worker's side and the corresponding types.}
\label{tab:labor-phi-mapping}
\end{center}

\textbf{No new family is needed.} A worker holds a composite contract combining families 1
and 5.

\begin{remark}[Not a special case but an object under different constraints]
\label{rem:worker-constraints}
Table~\ref{tab:labor-phi-mapping} shows that the catalogue applies once the viewpoint is
reversed, but
a worker cannot be regarded as a special case of a solo business. The constraints differ in
both $\Phi$ and $\phi$.

\textbf{Constraints on $\Phi$.} Five of the six degrees of freedom of Section~\ref{sec:dof}
are constrained by law. The Labour Standards Act requires wages to be paid in currency,
directly, in full, at least once a month, and on a fixed date; timing is fixed to arrears
and the paying party to the employer. State dependence is also permitted only to a limited
degree, so tying wages to performance is restricted.
\textbf{It is not that there is no choice, but that the choice set is small.}

\textbf{Constraints on $\phi$.} Defining $\phi_{\mathrm{prod}}$ for a worker would require a
marginal cost, and none is determined. For $\phi_{\mathrm{cog}}$ no counterpart is even
clear. The three-way decomposition of Chapter~\ref{sec:surplus} does not apply to workers as
it stands.

Furthermore the insolvency time of Section~\ref{sec:objective} is absorbing, whereas a worker
does not cease to exist upon being unable to pay. This text does not treat workers as an
independent object.
\end{remark}

Two points are worth recording. First, because monthly pay in arrears is the norm,
$\kappa_L>0$ dominates. Payroll lags are usually 15--45 days, so workers extend credit to
firms on the same order as trade payables. That $\kappa_L$ belongs in the third term
of~\eqref{eq:cash} should be made explicit.

Second, an employment contract is also a trade in which the stability of $\kappa$ is bought
at the price of assigning credit to the organization (Section~\ref{sec:credit-formation}).
In the language of the catalogue:
\begin{quote}
The worker avoids 5-1 (success fee, high variance) and chooses 1-3 (one-off arrears, low
variance), giving up the accumulation of public credit in exchange.
\end{quote}
This has the same form as Proposition~\ref{prop:discount}: whoever chooses stability
discounts a future claim and gives it away.

\subsection{Composite contracts: combinations of existing types}
\label{sec:composite}

Actual B2B contracts are often composites of several types. This is
\textbf{not a limit of the catalogue but the correct way to use it}.
By Proposition~\ref{prop:kappa-additive}, $\kappa$ is additive over parallel $\Phi$, so a
description that lists several types is meaningful at the level of $\kappa$.
But, as Remark~\ref{rem:decomposition-not-unique} states,
\textbf{the decomposition is not unique}. Several sets of types give the same $\kappa$, so
the types are a \textbf{generating set} for $\Phi$, not a basis.
The assignment procedure is fixed in Remark~\ref{rem:composite-assignment}.

\subsection{Claims that circulate: outside the framework}
\label{sec:token-limit}

Forms in which $\pi$ itself circulates as a claim on future $\delta$ --- token issuance, for
instance --- still cannot be handled.

The map $\Phi:\Delta\to\Pi$ of~\eqref{eq:phi} presumes that delivery and settlement
correspond between the same counterparty. When a claim is resold to a third party,
\textbf{the party that receives delivery changes after the fact}. This is a problem of the
domain and cannot be expressed in the setting of Chapter~\ref{sec:phi}, which fixes $N$.
An extension of the framework is required.

\section{A real firm is a bundle of \texorpdfstring{$\Phi$}{Phi}}

A large firm normally runs $\Phi$ from several families in parallel. Because the sign of
$\kappa$ differs by family, computing $\CCC$ at the level of the firm yields a mixed average
of terms with opposite signs. By Corollary~\ref{cor:ccc-weighted} that average is weighted
by the rate of sales, not a simple average of the families' $\CCC$.
To restate the conclusion of Section~\ref{sec:unit}:
\textbf{the unit of observation should be $\Phi$, not the firm}.

%============================================================
